\section{The bus and the routing}
\label{sec:bus}

Two nodes can be wired directly to each other. Five cannot --- that would be ten connections, and
the number grows quadratically. The solution is a \textbf{shared bus}: every node on the same
wire.

That has an immediate consequence: \textbf{everyone hears everything}. A bus is broadcast, and
that is not something you can opt out of. If you want a message to reach a particular node, you
have to write it in the message.

\subsection{Where the routing sits}

The bus's only job is to deliver every byte to every node. It does not read addresses, it knows
nothing about frames, it does not even know that a protocol exists.

\textbf{The routing sits in the node.} Every node receives every byte, feeds it to its parser,
and once a complete frame has been extracted compares the frame's \code{DST} with its own
address:

\begin{cppcode}
void Node::tick() noexcept
{
    if (!myParser.isFrameReady()) { return; }

    Frame frame{};
    if (!myParser.extractFrame(frame)) { myParser.reset(); return; }
    myParser.reset();

    if (frame.dstAddr != myAddr) { return; }   // not for us
    handleFrame(frame);
}
\end{cppcode}

Putting the routing in the node instead of in the bus is not a simplification made for the
course's sake. It is how a real bus works: a wire cannot read addresses.

\subsection{DST and SRC}

\code{DST} says who the frame is for. \code{SRC} says who it came from, and is needed so that the
receiver can answer --- without it, it does not know where the answer should go.

\code{SEQ} pairs a reply with its request. A node that has sent three requests and gets one reply
has to be able to tell which of them it answers.

\begin{aside}
\note CAN does all of this completely differently: there is no destination address at all, and
the identifier describes the \emph{content} of the message instead of its receiver. The comparison
is made in \kapref{ch:can}, and it is one of the most instructive in the course.
\end{aside}
